\documentclass[a4paper,12pt,twoside]{memoir}

% Castellano
\usepackage[spanish,es-tabla]{babel}
\selectlanguage{spanish}
\usepackage[utf8]{inputenc}
\usepackage[T1]{fontenc}
\usepackage{lmodern} % Scalable font
\usepackage{microtype}
\usepackage{placeins}
\usepackage{lscape}

\RequirePackage{booktabs}
\RequirePackage[table]{xcolor}
\RequirePackage{xtab}
\RequirePackage{multirow}

\usepackage{tcolorbox}

% Links
\usepackage[colorlinks]{hyperref}
\hypersetup{
	allcolors = {red}
}

% Ecuaciones
\usepackage{amsmath}

% Rutas de fichero / paquete
\newcommand{\ruta}[1]{{\sffamily #1}}

% Párrafos
\nonzeroparskip


% Imagenes
\usepackage{graphicx}
\newcommand{\imagen}[2]{
	\begin{figure}[!h]
		\centering
		\includegraphics[width=0.9\textwidth]{#1}
		\caption{#2}\label{fig:#1}
	\end{figure}
	\FloatBarrier
}

\newcommand{\imagenflotante}[2]{
	\begin{figure}%[!h]
		\centering
		\includegraphics[width=0.9\textwidth]{#1}
		\caption{#2}\label{fig:#1}
	\end{figure}
}



% El comando \figura nos permite insertar figuras comodamente, y utilizando
% siempre el mismo formato. Los parametros son:
% 1 -> Porcentaje del ancho de página que ocupará la figura (de 0 a 1)
% 2 --> Fichero de la imagen
% 3 --> Texto a pie de imagen
% 4 --> Etiqueta (label) para referencias
% 5 --> Opciones que queramos pasarle al \includegraphics
% 6 --> Opciones de posicionamiento a pasarle a \begin{figure}
\newcommand{\figuraConPosicion}[6]{%
  \setlength{\anchoFloat}{#1\textwidth}%
  \addtolength{\anchoFloat}{-4\fboxsep}%
  \setlength{\anchoFigura}{\anchoFloat}%
  \begin{figure}[#6]
    \begin{center}%
      \Ovalbox{%
        \begin{minipage}{\anchoFloat}%
          \begin{center}%
            \includegraphics[width=\anchoFigura,#5]{#2}%
            \caption{#3}%
            \label{#4}%
          \end{center}%
        \end{minipage}
      }%
    \end{center}%
  \end{figure}%
}

%
% Comando para incluir imágenes en formato apaisado (sin marco).
\newcommand{\figuraApaisadaSinMarco}[5]{%
  \begin{figure}%
    \begin{center}%
    \includegraphics[angle=90,height=#1\textheight,#5]{#2}%
    \caption{#3}%
    \label{#4}%
    \end{center}%
  \end{figure}%
}
% Para las tablas
\newcommand{\otoprule}{\midrule [\heavyrulewidth]}
%
% Nuevo comando para tablas pequeñas (menos de una página).
\newcommand{\tablaSmall}[5]{%
 \begin{table}
  \begin{center}
   \rowcolors {2}{gray!35}{}
   \begin{tabular}{#2}
    \toprule
    #4
    \otoprule
    #5
    \bottomrule
   \end{tabular}
   \caption{#1}
   \label{tabla:#3}
  \end{center}
 \end{table}
}

%
% Nuevo comando para tablas pequeñas (menos de una página).
\newcommand{\tablaSmallSinColores}[5]{%
 \begin{table}[H]
  \begin{center}
   \begin{tabular}{#2}
    \toprule
    #4
    \otoprule
    #5
    \bottomrule
   \end{tabular}
   \caption{#1}
   \label{tabla:#3}
  \end{center}
 \end{table}
}

\newcommand{\tablaApaisadaSmall}[5]{%
\begin{landscape}
  \begin{table}
   \begin{center}
    \rowcolors {2}{gray!35}{}
    \begin{tabular}{#2}
     \toprule
     #4
     \otoprule
     #5
     \bottomrule
    \end{tabular}
    \caption{#1}
    \label{tabla:#3}
   \end{center}
  \end{table}
\end{landscape}
}

%
% Nuevo comando para tablas grandes con cabecera y filas alternas coloreadas en gris.
\newcommand{\tabla}[6]{%
  \begin{center}
    \tablefirsthead{
      \toprule
      #5
      \otoprule
    }
    \tablehead{
      \multicolumn{#3}{l}{\small\sl continúa desde la página anterior}\\
      \toprule
      #5
      \otoprule
    }
    \tabletail{
      \hline
      \multicolumn{#3}{r}{\small\sl continúa en la página siguiente}\\
    }
    \tablelasttail{
      \hline
    }
    \bottomcaption{#1}
    \rowcolors {2}{gray!35}{}
    \begin{xtabular}{#2}
      #6
      \bottomrule
    \end{xtabular}
    \label{tabla:#4}
  \end{center}
}

%
% Nuevo comando para tablas grandes con cabecera.
\newcommand{\tablaSinColores}[6]{%
  \begin{center}
    \tablefirsthead{
      \toprule
      #5
      \otoprule
    }
    \tablehead{
      \multicolumn{#3}{l}{\small\sl continúa desde la página anterior}\\
      \toprule
      #5
      \otoprule
    }
    \tabletail{
      \hline
      \multicolumn{#3}{r}{\small\sl continúa en la página siguiente}\\
    }
    \tablelasttail{
      \hline
    }
    \bottomcaption{#1}
    \begin{xtabular}{#2}
      #6
      \bottomrule
    \end{xtabular}
    \label{tabla:#4}
  \end{center}
}

%
% Nuevo comando para tablas grandes sin cabecera.
\newcommand{\tablaSinCabecera}[5]{%
  \begin{center}
    \tablefirsthead{
      \toprule
    }
    \tablehead{
      \multicolumn{#3}{l}{\small\sl continúa desde la página anterior}\\
      \hline
    }
    \tabletail{
      \hline
      \multicolumn{#3}{r}{\small\sl continúa en la página siguiente}\\
    }
    \tablelasttail{
      \hline
    }
    \bottomcaption{#1}
  \begin{xtabular}{#2}
    #5
   \bottomrule
  \end{xtabular}
  \label{tabla:#4}
  \end{center}
}
%
% Nuevo comando para resaltar los términos en inglés.
\newcommand{\english}[1]{\emph{#1}}

%
% Nuevo comando para resaltar conceptos.
\newcommand{\concept}[1]{\textbf{#1}}

%
% Nuevo comando para resaltar nombres de archivo.
\newcommand{\filename}[1]{\texttt{#1}}

%
% Nuevo comando para URLs.
\newcommand{\myurl}[2]{\href{#1}{#2}\footnote{\url{#1}}}


\definecolor{cgoLight}{HTML}{EEEEEE}
\definecolor{cgoExtralight}{HTML}{FFFFFF}

%
% Nuevo comando para tablas grandes sin cabecera.
\newcommand{\tablaSinCabeceraConBandas}[5]{%
  \begin{center}
    \tablefirsthead{
      \toprule
    }
    \tablehead{
      \multicolumn{#3}{l}{\small\sl continúa desde la página anterior}\\
      \hline
    }
    \tabletail{
      \hline
      \multicolumn{#3}{r}{\small\sl continúa en la página siguiente}\\
    }
    \tablelasttail{
      \hline
    }
    \bottomcaption{#1}
    \rowcolors[]{1}{cgoExtralight}{cgoLight}

  \begin{xtabular}{#2}
    #5
   \bottomrule
  \end{xtabular}
  \label{tabla:#4}
  \end{center}
}


















\graphicspath{ {./img/} }

% Capítulos
\chapterstyle{bianchi}
\newcommand{\capitulo}[2]{
	\setcounter{chapter}{#1}
	\setcounter{section}{0}
	\chapter*{#2}
	\addcontentsline{toc}{chapter}{#2}
	\markboth{#2}{#2}
}

% Apéndices
\renewcommand{\appendixname}{Apéndice}
\renewcommand*\cftappendixname{\appendixname}

\newcommand{\apendice}[1]{
	%\renewcommand{\thechapter}{A}
	\chapter{#1}
}

\renewcommand*\cftappendixname{\appendixname\ }

% Formato de portada
\makeatletter
\usepackage{xcolor}
\newcommand{\tutor}[1]{\def\@tutor{#1}}
\newcommand{\course}[1]{\def\@course{#1}}
\definecolor{cpardoBox}{HTML}{E6E6FF}
\def\maketitle{
  \null
  \thispagestyle{empty}
  % Cabecera ----------------
\noindent\includegraphics[width=\textwidth]{cabecera}\vspace{1cm}%
  \vfill
  % Título proyecto y escudo informática ----------------
  \colorbox{cpardoBox}{%
    \begin{minipage}{.8\textwidth}
      \vspace{.5cm}\Large
      \begin{center}
      \textbf{TFG del Grado en Ingeniería Informática}\vspace{.6cm}\\
      \textbf{\LARGE\@title{}}
      \end{center}
      \vspace{.2cm}
    \end{minipage}

  }%
  \hfill\begin{minipage}{.20\textwidth}
    \includegraphics[width=\textwidth]{escudoInfor}
  \end{minipage}
  \vfill
  % Datos de alumno, curso y tutores ------------------
  \begin{center}%
  {%
    \noindent\LARGE
    Presentado por \@author{}\\ 
    en Universidad de Burgos --- \@date{}\\
    Tutor: \@tutor{}\\
  }%
  \end{center}%
  \null
  \cleardoublepage
  }
\makeatother

\newcommand{\nombre}{Adrián Arnaiz Rodríguez} %%% cambio de comando

% Datos de portada
\title{PDDetection\\ Aplicación de técnicas de minería de datos para la detección de la enfermedad del Parkinson}
\author{\nombre}
\tutor{José Francisco Díez Pastor y César Ignacio García Osorio}
\date{\today}

\begin{document}

\maketitle


\newpage\null\thispagestyle{empty}\newpage


%%%%%%%%%%%%%%%%%%%%%%%%%%%%%%%%%%%%%%%%%%%%%%%%%%%%%%%%%%%%%%%%%%%%%%%%%%%%%%%%%%%%%%%%
\thispagestyle{empty}


\noindent\includegraphics[width=\textwidth]{cabecera}\vspace{1cm}

\noindent D. José Francisco Díez Pastor y D. César Ignacio García Osorio, profesores del departamento de Ingeniería Civil, área Lenguajes y Sistemas Informáticos.

\noindent Exponen:

\noindent Que el alumno D. \nombre, con DNI 71306880A, ha realizado el Trabajo final de Grado en Ingeniería Informática titulado título de TFG. 

\noindent Y que dicho trabajo ha sido realizado por el alumno bajo la dirección del que suscribe, en virtud de lo cual se autoriza su presentación y defensa.

\begin{center} %\large
En Burgos, {\large \today}
\end{center}

\vfill\vfill\vfill

% Author and supervisor
\begin{minipage}{0.45\textwidth}
\begin{flushleft} %\large
Vº. Bº. del Tutor:\\[2cm]
Dr. D. José Francisco Díez Pastor
\end{flushleft}
\end{minipage}
\hfill
\begin{minipage}{0.45\textwidth}
\begin{flushleft} %\large
Vº. Bº. del co-tutor:\\[2cm]
Dr. D. César Ignacio García Osorio
\end{flushleft}
\end{minipage}
\hfill

\vfill

% para casos con solo un tutor comentar lo anterior
% y descomentar lo siguiente
%Vº. Bº. del Tutor:\\[2cm]
%D. nombre tutor


\newpage\null\thispagestyle{empty}\newpage




\frontmatter

% Abstract en castellano
\renewcommand*\abstractname{Resumen}
\begin{abstract}
La enfermedad del Parkinson causa, entre otras alteraciones, discapacidades en el habla a más del 90\% de los pacientes. Esta enfermedad causa trastornos del habla como pueden ser monotonía, volumen monótono, imprecisión en la articulación y otros síntomas. Hay numerosos estudios que abordan el problema de desarrollar aplicaciones de ayuda de evaluación del Parkinson.

En esta investigación, a partir de una base de datos de audios adquirida, tratamos de realizar todo el proceso de minería de datos que nos permita llegar a construir un clasificador. Este clasificador distingue entre si la persona de un audio tiene Parkinson o no. Para ello, se abordarán diferentes tipos de características extraídas de los audios (características físicas y más novedosas) con las que se realizarán diferentes experimentos.

Por ello, se realizará una aplicación de ayuda al facultativo médico que, usando el clasificador desarrollado en la investigación, sea capaz de agilizar el proceso de diagnosis de la enfermedad del Parkinson, además de reducir el riesgo de fallo en la misma.
\end{abstract}

\renewcommand*\abstractname{Descriptores}
\begin{abstract}
\textit{Minería de datos, Clasificador, Análisis del discurso, Parkinson, Aprendizaje automático, Aprendizaje profundo, Procesamiento de audio, Predicción enfermedades.}
\end{abstract}

\clearpage

% Abstract en inglés
\renewcommand*\abstractname{Abstract}
\begin{abstract}
Parkinson's Disease causes speech impairments in more than 90\% of patients, among other disabilities. This disease causes different disorders like monotony, monotonous volume, inaccuracy articulation and other symptoms. There are numerous studies that have addressed the problem of developing applications for helping the diagnosis of Parkinson's.

In this research, with a database of acquired audios, we try to realice the whole data mining process, building a good classifier as the main achieve. This classifier distinguishes between if the person of an audio has Parkinson's or not. For this purpose, we address different experiments, which are carried out with different types of characteristics extracted from the audios (physical characteristics and newer ones).

Therefore, we will make an aid application addressed to medical professional who, using the classifier developed in the research, will be able to speed up the diagnosis process of Parkinson's disease, in addition to reducing the risk of failure in it.
\end{abstract}

\renewcommand*\abstractname{Keywords}
\begin{abstract}
\textit{Data mining, Clasificator, Speech analysis, Parkinson, Machine learning, Deep learning, Audio processing, Disease prediction.}
\end{abstract}

\clearpage

% Indices
\tableofcontents

\clearpage

\listoffigures

\clearpage

\listoftables
\clearpage

\mainmatter
\capitulo{1}{Introducción}

\section{Introducción al proyecto}
En enfermedades como la depresión y las neurodegenerativas es requerido un seguimiento cercano de la enfermedad para ver su evolución. Actualmente es necesario que un facultativo médico se desplace a los domicilios de los pacientes y realice una serie de test y entrevistas estructuradas para evaluar la progresión de la enfermedad (depresión, alzheimer, \textbf{Parkinson}, etc). En líneas generales se pretende que el sistema propuesto sea capaz de determinar esa diagnosis solamente a partir de grabaciones de voz. En la actualidad, la monitorización de los síntomas y la diagnosis de enfermedades son costosos y, logísticamente, son inconveniente para el paciente y el personal clínico, lo que también dificulta el reclutamiento para futuros ensayos clínicos a gran escala. Este problema se acentúa en zonas rurales, donde se juntan otros factores que agravan el problema logístico: la falta de facultativos sanitarios, la dificultad de movilidad y la alta edad media de las personas residentes en los pueblos. Por ello, todo lo que signifique reducir complejidad y costes en este ámbito ayudará tanto a pacientes como a profesionales. Además, estas pruebas diagnósticas se realizan utilizando solamente pruebas de voz simples, auto-administradas y no invasivas.

El deterioro del habla del Parkinson se ha estudiado utilizando grabaciones de voz o conjuntos de datos de voz. Los conjuntos de datos de voz son características extraídas de las grabaciones de voz. Para ello, se pueden utilizar grabaciones de voz de diversos tipos, pero siempre de pronunciaciones estructuradas como monólogos, frases, palabras y vocales. El autor Orozco-Arroyave estudio estos 4 tipos en la investigación reflejada en el artículo \cite{Orz2016} (ver sección \ref{cap:TrabRel}). Sin embargo, las grabaciones de vocales sostenidas (pronunciación continua de una vocal durante un número determinado de segundos) son cada vez más populares como la entrada analítica. Holmes \cite{j2000voice}, estudió las características de la voz para etapas tempranas y posteriores de PD en comparación con personas sanas usando 4 segundos de vocal sostenida `a' y grabaciones de monólogos de 1 minuto. 
Holmes utilizó las grabaciones para extraer ocho características para el análisis acústico a partir de la mitad de la vocal sostenida y el canto a escala. Estas características incluyen la frecuencia fundamental y su desviación estándar, la intensidad media y su desviación estándar, el rango de frecuencia de frecuencia máxima, la fluctuación de fase, el brillo y la relación de ruido a armónicos. Por ejemplo, en este estudio Holmes descubrió, entre otros muchos hallazgos, que en comparación con los controles y los datos normativos publicados previamente, las voces de los pacientes con PD \footnote{Parkinson Disease} en etapa temprana y tardía se caracterizaron perceptualmente por la variabilidad de tono y volumen limitados, la transpirabilidad, la dureza y el volumen reducido.

Desde entonces, un número cada vez mayor de investigadores han estado utilizando la vocal sostenida `a' para estudiar la disfonía de la enfermedad del Parkinson, incluidos los investigadores Little y Tsanas 2010. En esta investigación realizada por Little y Tsanas en el 2010 \cite{MxLtAccurate}, demostraron la replicación rápida y remota de la evaluación UPDRS (escala de calificación de la enfermedad de Parkinson, véase sección \ref{sec:Updrs}) con una precisión clínicamente útil (aproximadamente 7.5 puntos UPDRS diferencia de las estimaciones de los clínicos).

Por lo tanto, nuestro proyecto consistirá en lo siguiente. Debemos conseguir una base de datos de audios. Estos audios se deberán realizar las entrevistas estructuradas, en las cuales sus respuestas serán grabadas en el mismo momento de la realización. Las grabaciones se identificarán con el identificador del paciente, así como con datos adicionales, como edad y sexo, del paciente. Las grabaciones incluirán grabaciones estandar, donde siempre se recite la misma frase, la misma palabra o la misma vocal. Concretamente, la base de datos con la que trabajaremos se describe en  \cite{OrzCorpus}. Los audios de las grabaciones se procesarán para obtener diferentes atributos y características de la voz, como por ejemplo los usados en el estudio de Orozco et al. \cite{Orz2016}, o los anteriormente comentados que se usaron en el estudio de Holmes. En nuestro caso se obtendrán los indicadores tono, cadencia, volumen y ritmo para su posterior análisis. A partir de ahí se construirá y entrenará una serie de modelos para predecir si la persona de un audio tiene Parkinson o no. Se realizarán una serie de diferentes experimentos con el objetivo de crear el mejor modelo posible. Este modelo correctamente construido se utilizará para evaluar la enfermedad de un paciente. Con ello conseguiremos avanzar en las investigaciones relativas a las enfermedades neurodegenerativas y su diagnóstico, y también realizar un producto, aplicación, que ayude e intente mejorar la calidad de vida tanto a pacientes como a trabajadores sanitarios.





\capitulo{2}{Objetivos del proyecto}

\section{Objetivos generales}
\begin{itemize}
	\item Investigar y condensar el estado del arte de la investigación sobre la detección del Parkison a través de la voz resumiendo los artículos científicos más importantes, identificando las tecnologías, herramientas y procesos actuales utilizadas en ellos, realizando taxonomías, etc.
	\item Recopilación de bases de datos adecuadas para la investigación, tanto para uso en este proyecto como para su uso en proyectos posteriores, cerciorando que son \english{datasets} de audios correctos para las tareas necesitadas (i.e. audios etiquetados).
	\item Realización de un estudio comparativo en cuanto a la utilización de diferentes modelos de clasificación y conjuntos de características extraídas de los audios. Se compararán resultados, tanto entre los diferentes experimentos que nosotros realicemos, como entre nuestros mejores experimentos y resultados científicos de anteriores experimentos (publicados en artículos científicos).
	\item Aportación de una nueva perspectiva a este campo de investigación: extracción de las características de los audios mediante \english{Deep Learning}.
	\item Finalmente, utilizar todo lo descrito anteriormente para realizar una aplicación la cual permita la monitorización de la diagnosis de la enfermedad del Parkinson a través de la voz. La aplicación será capaz de discernir, mediante un clasificador, si la persona de un audio subido a la aplicación web tiene Parkinson o no. Esta aplicación será un ejemplo de como poder plasmar los resultados de esta investigación en una herramienta funcional.
\end{itemize}

\section{Objetivos técnicos}
\begin{itemize}
	\item Desarrollar un algoritmo, cuya implementación en Python permita la extracción de características de los audios de manera adecuada (envolver la extracción de características que realizan Disvoice y VGGish en clases).
	\item Desarrollar una aplicación de escritorio para la etapa de explotación, que muestre como se pueden plasmar los resultados de la investigación para en una herramienta funcional. 
	\item Utilizar las herramientas más adecuadas para la realización del estudio comparativo: valoraremos Python y sus bibliotecas, Weka...
	\item Utilizar un sistema de control de versiones Git utilizando para ello la plataforma Github. Utilizar un cliente Git para el trabajo local para lo cual utilizaremos TortoiseGit.
	\item Utilizar la metodología Scrum en la elaboración del proyecto y concretamente la herramienta ZenHub (como extensión de Github) para la ayuda en la gestión de proyectos.
	
	
\end{itemize}
\capitulo{3}{Conceptos teóricos}

Este proyecto abordará diferentes temas dentro del campo conocido como \english{Data Mining}. A su vez se tratan temas sobre el análisis del discurso desde una perspectiva física. Se explicarán conceptos dentro de ambos campos para la correcta comprensión del proyecto.


\section{Minería de datos y aprendizaje automático}

La \concept{minería de datos} es un campo del conocimiento dentro de las ciencias de computación cuyo objetivo general es analizar grandes volúmenes de datos para extraer conocimiento de ellos. Se basa en el concepto de que la sociedad, sobre todo actualmente, produce una gran cantidad de datos y de diferente origen. Sin embargo, extraer conocimiento de los datos no es un proceso tan trivial. No es posible sacar la información de los datos en crudo. Por ello se necesitan diferentes métodos y técnicas que nos permitan hacer esa tarea, como son las técnicas de \concept{aprendizaje automático} \cite{datamining}. Este proceso de divide en varias fases.

\imagen{FasesDM}{Fases en la extracción de conocimiento de un conjunto de datos. \cite{DMprocess}}

Podemos hablar de que la minería de datos es el análisis de un conjunto grande de datos para la obtención de conocimiento de ellos, ya sea relaciones entre esos datos o patrones ocultos, que nos permita obtener una ventaja competitiva o un conocimiento novedoso. Esto está relacionado con el proyecto en cuanto a que queremos obtener conocimiento de la relación voz-enfermedad del Parkinson a través un gran conjunto de audios. A continuación, definiremos varios conceptos dentro de este campo que son manejados en el proyecto.

\subsection{Pre-procesamiento de datos}

Como hemos comentado anteriormente, los datos no pueden ser tratados de manera directa, necesitan ser tratados en una etapa anterior al descubrimiento de información: pre-procesamiento. Esto es indispensable en la fase de pre-procesamiento, transformar los datos en bruto a otro conjunto de datos que pueda servir para ser procesado por los algoritmos necesarios. Específicamente en nuestro caso, no podemos obtener conocimiento de el conjunto de audios en bruto, estos audios deberán ser tratados para extraer un conjunto de características numéricas de cada uno de ellos y poder seguir avanzando a partir de ese punto. El conjunto de características extraídas de los audios será en su mayoría propiedades físicas (\ref{sec:ccas}) o características extraídas con bibliotecas de \english{Deep Learning}. Como nuestro proyecto se basará en aprendizaje supervisado, también se llevará a cabo en esta etapa el etiquetado de los datos.

Otra parte importante, aunque no indispensable, dentro de esta etapa es la fase de selección de características. Cuando tenemos un gran número de atributos para cada instancia puede ser recomendable reducir su dimensionalidad. Una manera de reducir su dimensionalidad es a través de la selección de atributos manual, mediante un análisis estadístico \cite{DMprocess}. Un ejemplo de su uso puede verse en \cite{MxLtNovel}, donde varios de ellos son utilizados para definir un subconjunto de 10 características de cada audio dentro de las 132 medidas que se obtenían inicialmente en cada uno de ellos. En nuestro caso, no se realiza ese proceso manualmente. Por un lado, hemos realizado algunos experimentos donde hemos proporcionado todos los atributos de cada audio al algoritmo de clasificación y trasladando a éste la tarea de discernir cuales serán más importantes. Por otro, hemos realizado experimentos donde, en el proceso, se utilizan diferentes métodos selectores de atributos como \english{Select K Best} o \english{Variance Treshold} (ver sección \ref{subsec:selectores}). 

\subsection{Aprendizaje automático}
Es el campo de las ciencias de la computación enfocado a que los dispositivos \textit{aprendan} por ellos mismos sin haber sido explícitamente programados para ello. Esta es una rama de la inteligencia artificial. Tras obtener un conjunto válido de características se aplican algoritmos de aprendizaje automático. Estos algoritmos son los encargados de detectar patrones en los datos. Estos algoritmos tienen un fuerte componente matemático y estadístico ya que, a través de métodos de estos campos del conocimientos, se extrae la información de los datos. Se divide en:
\begin{itemize}
	\item \concept{Aprendizaje supervisado}: para cada instancia del conjunto de datos tendremos tanto la entrada como la salida deseada. El objetivo será predecir la salida correcta para una nueva entrada. En nuestro caso la entrada serían las características de los audios y la salida si la persona de ese audio tiene parkinson (clasificación) o en qué nivel lo tiene (regresión) \footnote{Estos ejemplos ilustran la diferencia entre clasificación, donde hay que predecir una etiqueta de entre un conjunto finito de etiquetas posibles, y regresión, donde el valor que se predice es una magnitud continua de un conjunto con un número de valores potencialmente infinito.}.
	\item \concept{Aprendizaje no supervisado}: únicamente se tienen los datos de entrada y no se sabe nada acerca de su salida o clase. Habitualmente estos algoritmos son utilizados para la detección de clases que permitan separar los datos en diferentes grupos. Nosotros no utilizaremos este tipo en nuestro proyecto.
\end{itemize}

\subsubsection{Generalización}
El objetivo de los algoritmos de aprendizaje supervisado no es lograr dar salidas correctas para los ejemplos de datos que tenemos, sino lograr clasificar o dar la salida correcta para un nuevo ejemplo que nos llegue. Se puede decir que el objetivo del aprendizaje automático es en realidad es obtener conocimiento a partir de un gran conjunto de datos el cual sea extrapolable para todos los datos dentro de esa misma situación (de ese mismo problema).
El problema de ajustarse demasiado a los datos de entrenamiento perdiendo así capacidad de generalización se llama sobreajuste. Esto ocurre cuando se crean modelos demasiado complejos que se ajustan a los datos de entrenamiento al 100\%. Al ocurrir sobreajuste, corremos el riesgo de clasificar mal nuevos datos. Si además este sobreajuste a los datos de entrenamiento se combina con la aparición de datos ruidosos en nuestro dataset, se perderá aún más capacidad de generalización. Por ello es importante tener métodos de evaluación de modelos que nos permitan medir cuánto de bueno es nuestro modelo a la hora de generalizar. Aquí aparece el concepto \concept{Cross-Validation}.

\subsubsection{$K$-fold Cross Validation \cite{CrossVal}}
La validación cruzada es un método para la evaluación y comparación de algoritmos. Se basa en descomponer el conjunto de datos en 2 subconjuntos, entrenamiento y test, con el objetivo de evaluar que tal generaliza nuestro modelo. Sin embargo, de esta manera estamos perdiendo información de nuestros datos, ya que el conjunto que utilizamos para test nunca es usado para entrenar. La solución a este problema es la validación cruzada $k$-fold.

En \concept{k-fold cross-validation} el conjunto entero de datos se divide en $k$ conjuntos del mismo tamaño cada uno. A partir de ahí se itera sobre los subconjuntos $k$ veces utilizando cada vez $k-1$ conjuntos para entrenar y 1 conjunto de test. Para cada una de las iteraciones se guarda la performance y finalmente se hace la media para obtener una performance total. El número idóneo de $k$ es 10 según \cite{CrossVal} y así es utilizado el 10-fold en \cite{Orz2016}.

\imagen{3fold_cv}{Ejemplo 3-fold cross-validation extraído de \cite{CrossVal}.}

Cabe destacar que \cite{CrossVal} explica que para \english{k-fold cross-validation} se debe realizar estratificación. \concept{Estratificación} es el proceso por el cual se asegura que una de las $k$ particiones de los datos total es una buena representación de los datos totales. Por ejemplo, en nuestro caso tenemos 50\% de personas con Parkinson y 50\% de personas sanas. Esto significa que para estratificar de manera correcta, cada partición deberá tener la mitad de instancias de cada clase. Cada partición $k$ deberá tener aproximadamente la misma proporción de clases que el conjunto total.

\subsection{Nested Cross Validation}
A menudo, queremos ajustar los parámetros de un modelo (por ejemplo, $C$ en SVM). Es decir, queremos encontrar el valor de un parámetro que maximice el rendimiento de nuestro clasificador. Sin embargo, Cawley y Talbot en \cite{cawley2010over}, aseguran que si utilizamos la partición de test de los datos, tanto para ajustar los parámetros, como para evaluar el modelo, corremos el riesgo de sesgar de forma optimista las evaluaciones del modelo. Por esta razón, si un conjunto de pruebas se utiliza para seleccionar parámetros del modelo, entonces necesitamos un conjunto de pruebas diferente para obtener una evaluación imparcial de ese modelo seleccionado.


Una forma de superar este problema es tener validaciones cruzadas anidadas (\english{Nested Cross Validation}). Primero, se usa una validación cruzada interna para ajustar los parámetros y seleccionar el mejor modelo (o parámetros del modelo). En segundo lugar, se utiliza una validación de cruz externa para evaluar el modelo seleccionado por la validación de cruz interna. Ver Figura \ref{fig:nested-k-fold}.

\imagen{nested-k-fold}{Validación cruzada anidada}



\subsection{Algoritmos de aprendizaje automático}

A continuación explicaremos varios algoritmos de los utilizados en este proyecto.

\subsubsection{SVM \cite{svmexpl}}

Este algoritmo de aprendizaje supervisado es llamado como máquinas de vector de soporte (SVM). Es un clasificador lineal basado en el concepto de margen máximo. Se utiliza tanto para regresión como para clasificación. Para clasificación su principio fundamental es separar dos clases en el conjunto de datos utilizando un hiperplano. El hiperplano utilizado como separador de las clases será el que maximice el margen entre las clases. Como sabemos, los problemas del mundo real tienen complicaciones como que no sean linealmente separables, que las clases estén solapadas, que haya más de 2 clases de datos, etc. 

Para solucionar el problema de la existencia de más de 2 clases, se utiliza un SVM por cada una de las clases. Sin embargo, para el problemas de separar dos clases que no son linealmente separables, utilizamos el \concept{kernel trick}. Esta característica del algoritmo es importante para poder separar este tipo de datos. Estas funciones kernel son utilizadas para aumentar la dimensionalidad de los datos de entrada, es decir, transformar el espacio de entrada. De esta manera dividiendo linealmente el espacio transformado con el hiperplano pueden resolverse el problema de la separación lineal.  Existen diferentes ejemplos de funciones kernel, entre las más habituales están el kernel, polinomial, funciones de base radial, sigmoide y kernel Gaussiano que es utilizado en \cite{MxLtNovel} y \cite{Orz2016}.

\subsubsection{Random Forest \cite{breiman2001random}}

Es un método de \concept{ensemble} (métodos combinados) basado en los árboles de predicción, la combinación de la selección aleatoria de atributos y el \english{bagging}. El proceso es el siguiente: se construirán de manera independiente un conjunto de árboles, los cuales serán todos diferentes, debido a que están construidos con \english{bagging} y aleatoriedad de atributos, a la hora de predecir una nueva instancia, cada uno hace su predicción y, teniendo en cuenta la decisión de todos los árboles, llegamos a una decisión final conjunta. \concept{Bagging} consiste en obtener varios subconjuntos de datos a partir de un único conjunto mediante remuestreo con reemplazamiento. Esto se hace escogiendo $N$ elementos del conjunto inicial de manera aleatoria pudiendo coger el mismo ejemplo varias veces. Esto se acopla a la construcción de cada árbol de la siguiente manera:
\begin{enumerate}
	\item A la hora escoger el atributo discriminante de un nodo, el mejor lo elegiremos de entre un subconjunto de todos los atributos. El mejor tamaño de este subconjunto es de tamaño $\sqrt{M}$ o $\log M$, siendo $M$ el número de atributos.
	\item Se realiza ese proceso para cada nodo del árbol.
\end{enumerate}
A la hora de realizar la predicción del error se sigue el método \concept{out of bag}. Consiste en recorrer las instancias de entrenamiento y predecir la clase de cada una únicamente con el conjunto de árboles que no han tenido esa instancia en su conjunto de entrenamiento. Se hace ese proceso para todas las instancias que tenemos y predecimos el error.

\subsection{Algoritmos de selección de atributos} \label{subsec:selectores}
Cuando nosotros tenemos un conjunto de datos, a veces es necesario saber qué atributos de los que tenemos con los más importantes. Puede ocurrir que haya atributos irrelevantes, lo que puede producir sobreajuste. También puede haber atributos redundantes, lo que para ciertos algoritmos de aprendizaje es nocivo. O, simplemente, que tenemos una dimensión de datos de entrada muy grande, es decir, muchos atributos para cada instancia que, juntado al hecho de tener pocas instancias, puede ser un problema. Estos algoritmos seleccionan atributos en función de diferentes análisis: analizando cada atributo individual, utilizando un modelo de aprendizaje o evaluando grupos de atributos. A continuación, explicaremos los dos algoritmos de selección de atributos más usados en los experimentos del proyecto.

\subsubsection{Variance Treshold}
Este algoritmo selecciona atributos analizándolos de manera individual. Este algoritmo de selección de atributos se basa en eliminar atributos de baja varianza. Es decir, los atributos con poca varianza, aquellos que son casi idénticos en todos los ejemplos, son eliminados. Esto se realiza de esta manera, debido a que, los atributos con poca varianza, no aportan mucha información al ser muy parecidos en todos los ejemplos, lo que se traduce en que no nos aporta información para discernir si es de una clase u otra.

\subsubsection{Seleck K best}
Este algoritmo también selecciona atributos analizándolos de manera individual. Se basa en que a veces existen atributos poco discriminativos con la clase (el valor del atributo varía sin ninguna relación con la clase). En este método de selección, se aplica un test estadístico entre cada uno de los atributos y la clase, quedándonos con los $k$ mejores, siendo $k$ un parámetro ajustable. 

\subsection{Optimización de parámetros} 
Una de las tareas más complejas en la minería de datos, es acertar con la parametrización de los modelos. Para un mismo tipo de modelo, puede variar mucho el rendimiento dependiendo el valor de sus parámetros. Pero la tarea de saber que parámetros son los mejores no es una tarea trivial. Requiere de mucho conocimiento del propio modelo, circunstancia que la mayoría de las veces no ocurre.

Para ello, una de las soluciones que podemos tener es probar muchos tipos de parámetros para elegir el mejor. En \textit{Python} esto se llama \textit{Grid Search}. Básicamente, se trata de definir los valores diferentes que queremos probar de cada parámetro o parámetros y probar todas las posibles combinaciones en el modelo. En la Figura \ref{fig:GridSearch}, vemos la importancia de la optimización de parámetros, donde se obtiene un 0.076 AUC más utilizando \textit{Grid Search}.

\imagen{GridSearch}{Importancia optimización de parámetros}

\section{Deep Learning}
El \textit{Deep Learning} o aprendizaje profundo, consiste en llevar a cabo el proceso de aprendizaje automático mediante redes neuronales de varias capas, en las que cada capa se encarga de extraer características a partir de las características de las capas anteriores. Para ello, se modelan arquitecturas complejas que consiguen transformaciones no lineales múltiples, a través de datos expresados de forma matricial. Dependiendo del número de capas, número de neuronas, funciones de activación y demás elementos de la red neuronal, se obtienen unas arquitecturas u otras. Cabe destacar que estos modelos funcionan como una caja negra y son poco interpretables.

\section{Conceptos estadísticos}
Se describirán a continuación una serie de conceptos matemático-estadísticos utilizados tanto en la evaluación de clasificadores, como en la obtención de medidas de los audios.
\subsection{ROC}
Un gráfico de características operativas del receptor (ROC) \cite{roc}, es una técnica para visualizar, organizar y seleccionar clasificadores según su rendimiento. Los gráficos ROC se han utilizado durante mucho tiempo en la teoría de detección de señales para representar la relación entre las tasas de aciertos y las falsas alarmas de los clasificadores. El análisis de ROC se ha extendido para su uso en la visualización y análisis del comportamiento de los sistemas de diagnóstico, como es el de nuestro caso. El gráfico ROC traza una curva de probabilidad. Visto de otra manera, puede representar el ratio de verdaderos positivos frente a la razón o ratio de falsos positivos. Es decir, hasta qué punto un clasificador es capaz de clasificar los casos positivos correctamente.

La medida de rendimiento de clasificadores que se saca de este gráfico es la llamada \english{Area Under Curve} (AUC). En líneas generales, esta medida indica \textbf{cómo de bien el modelo es capaz de distinguir entre clases}.  Un modelo que es capaz de separar perfectamente todas las clases tendrá un AUC cercano a 1, mientras que si es incapaz de separarlas, tendrá un AUC de 0,5. 




\subsection{Medidas de distribución}\label{subsec:medidasdistribucion}
Son diferentes medidas extraídas de cada medida física de los audios. Por ejemplo, de la amplitud de onda de un audio se sacan sus 4 funcionales correspondientes con sus momentos de distribución (media, desviación típica, coeficiente de asimetría, curtosis). Se obtienen debido a que caracterizan una muestra de tal manera que si dos distribuciones tienen los momentos iguales son iguales.
\begin{itemize}
	\item \concept{Media}: Es el resultado de la suma de todos los valores dividida entre el número de ellos. Se corresponde con el momento de la distribución de orden 1 respecto a la origen. $m$.
	\item \concept{Desviación}: Mide la dispersión que tienen unos datos respecto a la media. Cuanto mayor sea este valor significará que los datos se encuentran en un rango amplio respecto a la media, mientras que si su valor en bajo significará que los valores se agrupan en un rango cercano a la media. Se corresponde con el momento de la distribución de orden 2 respecto a la media. $d$.
	\item \concept{Oblicuidad o coeficiente de asimetría}: Mide la mayor o menor simetría de la distribución. Contra mayor sea el coeficiente de oblicuidad, mayor será simetría de los datos respecto a la media. Se corresponde con el momento de la distribución de orden 3 respecto a la media. $sk$. 
	\item \concept{Curtosis}: Mide la mayor o menor concentración de datos alrededor de la media. A un mayor valor de este coeficiente, se entiende que los valores están más agrupados en torno a la media y en valores alejados de ella, dejando los tramos intermedios con menor frecuencia. Se corresponde con el momento de la distribución de orden 4 respecto a la media. $k$.
\end{itemize}
Momento de la distribución de orden \textit{r} respecto a la media:\\
\begin{equation}
m_{r} = \sum_{i=1}^{n} (x_{i} - \overline{x})^{r} P(n_{i})
\end{equation}

\section{UPDRS}\label{sec:Updrs}
\english{Unified Parkinson's disease rating scale} \cite{updrs}, \concept{UPDRS} es la escala universal con la que se mide el grado de severidad del Parkinson. Fue creada para dar un estándar a la evaluación de áreas específicas de la discapacidad. Esta escala evalúa 6 partes fundamentales: 1-Comportamiento, 2-Actividades de la vida diaria, \textbf{3-Examen motor} \cite{updrs3}; 4-Complicaciones de la terapia, 5-Escala de Hoehn \& Yahr (Severidad) \cite{hoehn1967} y 6-Escala de Schwab y England (vida cotidiana).\\
\subsection{UPDRS-III}
Debido al objetivo de nuestro estudio, nos centraremos en la parte 3: examen motor \cite{updrs3}. Utiliza una escala de 0-4, donde 0 es la ausencia de discapacidad motora y 4 discapacidad motora severa. Esta parte trata diversos temas todos relacionados con la capacidad motora corporal. Mide entre otras cosas elementos como el habla, la expresión facial, el temblor o la agilidad. Todo esto afecta a la voz en cuanto se debe realizar un control correcto de los músculos glotales, labios, faringe y más para conseguir una pronunciación satisfactoria. La discapacidad motora  produce a la hora de hablar volúmenes bajos, discurso monótono, articulación imprecisa \cite{Orz2016}. El englobe de estas capacidades se denomina disartria hipocinética \cite{disartia}.

\subsection{Hoehn \& Yahr}
La escala de Hoehn \& Yahr \cite{hoehn1967} evalúa la enfermedad del Parkinson en una escala del 1 al 5. La enfermedad es más severa en función del aumento de la escala. El grado mínimo 1 se corresponde con que la enfermedad es exclusivamente unilateral, aumentando hasta el grado 5 en el que el paciente está en silla de ruedas o en cama, si no tiene ayuda.

\section{Análisis del discurso}
El análisis del discurso se realiza desde 3 prismas diferentes, que son la fonación, la articulación y la prosodia \cite{speechAnalysis}. En nuestro proyecto, fijaremos estos grupos de características a extraer de cada tipo de audios.

\subsection{Análisis de la fonación}
La fonación aborda la vibración de las cuerdas vocales a la hora producir un sonido \cite{speechAnalysis}. Desde el punto de vista clínico, este análisis está relacionado con la curvatura y el cierre incompleto de las cuerdas vocales a la hora de emitir un sonido \cite{phonationfeat}. Las medidas más típicas para el análisis de fonación son Jitter, Shimmer, APQ y PPQ \cite{neurospeech}. Explicadas en la sección \ref{sec:ccas}.

\subsection{Análisis de la articulación}
La articulación comprende la modificación de la posición, el estrés y la forma de los órganos y tejidos involucrados en la producción del habla \cite{speechAnalysis}. Esto se manifiesta en déficits como, por ejemplo, una reducción de la amplitud y la velocidad de los movimientos articulatorios de los labios, la mandíbula y la lengua \cite{articulationfeat}. La capacidad de articulación es evaluada a través de la energía que se libera en las transiciones entre segmentos \textit{voiced} $\rightarrow$ \textit{unvoiced} (transición entre segmentos con voz y sin voz) y viceversa con medidas como MFCC o BBE, explicadas en \ref{sec:ccas}. Se basa en que los pacientes de PD tienen dificultad para comenzar y para detener la vibración de las cuerdas vocales \cite{neurospeech}.

\subsection{Análisis de la prosodia}
La prosodia aborda temas como la variación del volumen, el tono y la sincronización para hablar de manera natural \cite{speechAnalysis}. Se manifiesta con monotonía en el volumen y en el tono, cambios de rapidez en el habla y dificultades a la hora de expresar emociones a través del discurso \cite{emotionfeat}. Las medidas más típicas para el análisis de prosodia son las relacionadas con la frecuencia fundamental, el contorno de la energía y la duración. Explicadas en la sección \ref{sec:ccas}.

\section{Características Físicas de la voz} \label{sec:ccas}
Desde el punto de vista ingenieril y relacionado con el aprendizaje automático, todos los análisis anteriores se deben expresar de una manera numérica para poder crear vectores de características de un audio. Por ello, se extraen diferentes características numéricas de los audios que intentan expresar los análisis de fonación, articulación y prosodia de la manera más óptima posible. Otra dificultad importante desde el punto de vista ingenieril es la continuidad de los audios. Los audios cambian de manera continua en el tiempo, lo que complica el proceso de la extracción de características. Por ello, el audio se divide pequeños segmentos de tiempo en los que se calculará las características de manera estática para ese tramo. Las medidas físicas extraídas de los archivos de voz son de gran variedad y por ello en este apartado explicamos las más comunes y más usadas en los diferentes estudios del análisis del discurso. 

\subsection{Jitter}
El concepto \concept{Jitter} es usado para evaluar la variabilidad temporal durante el envío de señales digitales \cite{wiki:jitter}. Concretamente, en el análisis del discurso, esta medida significa la variación temporal de la frecuencia fundamental del discurso \cite{neurospeech}. $N$ es el número de fragmentos, $M_{f}$ la frecuencia máxima y $F_{0}(k)$ la amplitud de ese frame en concreto:
\begin{equation}
\mathit{Jitter}(\%) = \frac{100}{N M_{f}} \sum_{k=1}^{N}|F_{0}(k)-M_{f}|
\end{equation}

\subsection{Shimmer}
El concepto \textbf{Shimmer} es usado, en el análisis del discurso, para medir la variación temporal de la amplitud del discurso \cite{neurospeech}. $N$ es el número de \english{frames}, $M_{a}$ la amplitud máxima y $A(k)$ la amplitud de ese \english{frame} en concreto:
\begin{equation}
\mathit{Shimmer}(\%) = \frac{100}{N M_{a}} \sum_{k=1}^{N}|A(k)-M_{a}|
\end{equation}

\subsection{APQ y PPQ}
\textbf{APQ} mide la variabilidad a largo plazo de la amplitud de la voz. Para calcularla, se suaviza mediante una media móvil de tamaño 11 y se calcula como la diferencia media absoluta entre la amplitud de un \english{frame} y las amplitudes promediadas sobre sus vecinos, dividida por la amplitud media.

La \textbf{PPQ} mide la variabilidad a largo plazo de la frecuencia fundamental y para calcularla se suaviza mediante una media móvil de tamaño 5. Se calcula como la diferencia promedio absoluta entre la frecuencia de cada cuadro y el promedio de sus vecinos, dividida por la frecuencia media. Ambas medidas se calculan de igual manera, APQ relativa a la amplitud y PPQ relativa a la frecuencia \cite{neurospeech}.

\subsection{Mel Frequency Cepstral Coefficients - MFCC}\label{subsec:mfcc}
Los \textbf{Coeficientes Cepstrales en las Frecuencias de Mel} son 12 coeficientes para la representación del habla basados en la percepción auditiva humana. Con ellos podemos obtener la información más relevante de una porción de audio, obviando partes menos importantes como el ruido. El cálculo de estos coeficientes está basado en la transformada de Fourier y la transformada del coseno discreta \cite{wiki:mfcc}. Como hemos comentado anteriormente, nos sirve para medir la energía que se libera en las transiciones entre segmentos \textit{voiced} $\rightarrow$ \textit{unvoiced} y viceversa. Tiene un crecimiento logarítmico y se calcula según la siguiente función:
\begin{equation}
m = 1127,01048\ln{\frac{1+f}{700}}
\end{equation}

\subsection{Bark Band Scale}
\textbf{Bark Band Scale} es otro método de caracterización de la energía liberada en las transiciones de la voz. Da 24 coeficientes acorde a la escala de Bark \cite{bbe}. Es una escala psicoacústica en la que se definen límites de frecuencias para definir la escala (i.e. grado 1 [100Hz-200Hz], grado 2 [200Hz-300Hz],..., grado 23 [12000-15500], grado 24 [15500- $\infty$])
Tiene un crecimiento logarítmico y se calcula según la siguiente función:
\begin{equation}
\mathit{Bark}(f) = 13 \arctan(0.00076 f)+ 3.5 \arctan \left ( \frac{f}{7500} \right )^{2}
\end{equation}


\subsection{Medidas relacionadas con la frecuencia fundamental}
Su objetivo es caracterizar los patrones de entonación de la voz. Por ello, se calcula el contorno de la frecuencia máxima y se calculan diferentes estadísticos sobre ella, como medias en Hz, desviación en Hz, máximos, etc \cite{neurospeech}.

\subsection{Medidas relacionadas con la energía}
Al igual que para el anterior caso, pero esta vez con la energía. Se calculan medidas estadísticas sobre el contorno de la energía: media en dB, desviación en dB, máximo o el coeficiente de regresión entre el contorno de energía y una regresión lineal \cite{neurospeech}.

\capitulo{4}{Técnicas y herramientas}

Explicaremos las herramientas utilizadas en nuestro proyecto.

\section{Python}
\textcolor{red}{DEFINIR}.

\section{Anaconda Distribution}
\textcolor{red}{DEFINIR}.

\section{Jupyter Notebook IDE}
\textcolor{red}{DEFINIR}.

\section{Librerías de Python}
\subsection{Numpy}
\textcolor{red}{DEFINIR}.

\subsection{Scikit-learn}
\textcolor{red}{DEFINIR}.

\subsection{Scipy}
\textcolor{red}{DEFINIR}.

\subsection{Pandas}
\textcolor{red}{DEFINIR}.

\subsection{Pysptk}
\textcolor{red}{DEFINIR}.

\subsection{Sounddevice}
\textcolor{red}{DEFINIR}.

\subsection{Pydub}
\textcolor{red}{DEFINIR}.

\subsection{Disvoice}
La biblioteca  \myurl{https://github.com/jcvasquezc/DisVoice}{Disvoice} \cite{neurospeech} es un conjunto de \english{scripts} de Python para la extracción de medidas del habla. Disvoice calcula medidas de articulación, de la fonación y de prosodia a partir de vocales sostenidas y expresiones verbales continuas, con el objetivo de evaluar las capacidades de comunicación de los pacientes con diferentes trastornos de la voz o trastornos neurodegenerativos como la enfermedad de Parkinson. Ha sido desarrolada por Juan Camilo Vásquez-Correa, el cual es co-autor de varios artículos con Juan Rafael Orozco-Arroyave como \cite{neurospeech}, y tiene licencia de software MIT. Cabe destacar que se han intercambiado dos correos electrónicos con JC Vásquez-Correa, en los que se nos explica tanto la utilización de los \english{scripts}, como la salida detallada de cada uno de ellos.

Contiene 3 \english{scripts} principales para la extracción de características de audios (fonación, articulación y prosodia), extrayendo por ejemplo de un audio hasta 488 medidas relacionadas con la articulación. Tiene una gran parte de su contenido dedicado a la visualización de los audios mediante diferentes métodos, característica que no es usada en nuestro proyecto. 

Esta biblioteca ha sido usada, como puede entenderse, para la etapa de extracción de características de los audios. Está escrita para ser usada como \english{scripts} de Python y utiliza internamente varias bibliotecas como scipy, numpy, scikitlearn, pysptk, sounddevice, os y programas como Praat \cite{praat}. A la hora de ser utilizada por nosotros ha tenido que ser configurada para su correcto funcionamiento en nuestro entorno. Los detalles de la configuración y los cambios realizados se dan en el manual del programador.


\subsection{Os}
\textcolor{red}{DEFINIR}.

\subsection{Keras}
\textcolor{red}{DEFINIR}. 

\subsubsection{VGGish}
\textcolor{red}{DEFINIR}. 

\subsubsection{VGGish2Keras}
\textcolor{red}{DEFINIR}.

\section{Praat}
\myurl{http://www.fon.hum.uva.nl/praat/}{Praat} \cite{praat} es un programa el cual nos permite realizar análisis fonéticos de audios vocales. Esta herramienta está enfocada a la investigación del habla. Permite hacer una multitud de análisis diferentes entre los que se encuentran análisis del discurso (análisis espectrales, análisis de intensidad, de formantes...) o análisis estadístico. Una característica importante para nosotros es que permite ser ejecutado mediante línea de comandos con diferentes parámetros. Otro aspecto interesante es que Praat tiene un \english{wrapper} para Python llamado Parselmouth \cite{parselmouth}, aunque nosotros no lo utilizamos en el proyecto. Ha sido desarrolada por Paul Boersma y David Weenink de la Universidad de Ámsterdam.

Será utilizada internamente por la biblioteca Disvoice para analizar una serie de características de los audios, que posteriormente volverá a procesar Disvoice con diferentes métodos de Python para devolvernos a nosotros las características finales deseadas.


\section{Git}
\textcolor{red}{DEFINIR}.

\section{Github}
\textcolor{red}{DEFINIR}.

\section{ZenHub}
\textcolor{red}{DEFINIR}.

\section{TortoiseGit}
\textcolor{red}{DEFINIR}.

\section{Latex}
\textcolor{red}{DEFINIR}.

\subsection{MiKTeK}
\textcolor{red}{DEFINIR}.

\subsection{\LaTeX}
\textcolor{red}{DEFINIR}.
\capitulo{5}{Aspectos relevantes del desarrollo del proyecto}

Este apartado pretende recoger los aspectos más interesantes del desarrollo del proyecto, comentados por los autores del mismo.
Debe incluir desde la exposición del ciclo de vida utilizado, hasta los detalles de mayor relevancia de las fases de análisis, diseño e implementación.
Se busca que no sea una mera operación de copiar y pegar diagramas y extractos del código fuente, sino que realmente se justifiquen los caminos de solución que se han tomado, especialmente aquellos que no sean triviales.
Puede ser el lugar más adecuado para documentar los aspectos más interesantes del diseño y de la implementación, con un mayor hincapié en aspectos tales como el tipo de arquitectura elegido, los índices de las tablas de la base de datos, normalización y desnormalización, distribución en ficheros3, reglas de negocio dentro de las bases de datos (EDVHV GH GDWRV DFWLYDV), aspectos de desarrollo relacionados con el WWW...
Este apartado, debe convertirse en el resumen de la experiencia práctica del proyecto, y por sí mismo justifica que la memoria se convierta en un documento útil, fuente de referencia para los autores, los tutores y futuros alumnos.

\capitulo{6}{Trabajos relacionados}\label{cap:TrabRel}

En este apartado se comentará el estado del arte de la materia y algunos trabajos relacionados.\\

Para comenzar, comentaremos que hay dos grupos de investigación importantes en este campo de estudio. El primer grupo sería el encabezado por el autor M. A. Litlle de la Universidad de Aston, Birmingham, Inglaterra. Este autor a su vez encabeza un proyecto para la obtención de una gran base de datos de 10000 audios obtenidos de pacientes con Parkinson. Esta iniciativa es conocida como \myurl{http://www.parkinsonsvoice.org/}{Parkinson Voice Inititative}. Este conjunto de autores tiene varios artículos \cite{MxLtAccurate}, \cite{MxLtNovel} y el más citado \cite{MxLtSuitability}. Estos artículos aportan un buen análisis sobre las características a extraer de los audios (lineares y no lineares) y algoritmos de selección de características. Los aspectos más relevantes de cada artículo son los siguientes:
\begin{itemize}
	\item \textit{Suitability of dysphonia measurements for telemonitoring of Parkinson’s disease} \cite{MxLtSuitability}.\\
	Este artículo trata sobre la aplicación de diferentes medidas estándares (lineales) y no estándares (no lineales) de disfonía para la clasificación automática de personas con Parkinson y personas sanas. Además, añaden otra nueva medida, calculada en este mismo artículo, llamada PPE, \english{pitch period entropy}. Se centra en responder a la pregunta de qué características son las mejores para la detección del Parkinson. Para ello utiliza únicamente audios de vocales sostenidas (28 PD y 8 HC). Se calculan un total de 17 medidas de disfonía de cada audio usando diferentes software como Praat \cite{praat}. Posteriormente, se hace un análisis de correlación entre las medidas obtenidas, y de aquellos pares que tienen un coeficiente de correlación mayor del 95\%, se elimina una. Después de este análisis se obtienen 10 características: \english{jitter} medido diferencia absoluta, \english{jitter} medido como media entre ciclos, APQ, \english{shimmer} (calculado como la diferencia absoluta promedio entre las amplitudes de los períodos consecutivos), HNR, NHR, RPDE, DFA, la dimensión de correlación y el comentado PPE.
	
	 En el siguiente paso del proceso, utilizan el algoritmo de aprendizaje supervisado SVM con kernel de base radial, para construir el modelo. Esto se realiza con cada uno de los 1023 diferentes subconjuntos posibles de las 10 características, $\sum_{i=1}^{10} \binom{10}{i}$, para encontrar el mejor subconjunto posible. Esto es debido a que los autores consideran que es un número pequeño de conjuntos y se pude hacer búsqueda exhaustiva. Se llega a la conclusión que el \textbf{mejor conjunto de características es el formado por HNR, RPDE, DFA, y PPE} que devuelve una precisión del \textbf{91\%}, seguido en precisión por el conjunto completo de las 10 características.
	
	\item \textit{Accurate Telemonitoring of Parkinsons Disease Progression by Noninvasive Speech Tests} \cite{MxLtAccurate}.\\
	En este artículo trata sobre cómo \textbf{el objetivo de la clasificación del Parkinson puede ser de regresión}. En vez de clasificar entre personas sanas (HC) y con Parkinson (PD), lo que se hace es medir el nivel de parkinson de una persona usando la escala UPDRS (\english{Unified Parkinson's Disease Rating Scale}). Para ello saca un total de 16 características de audios con vocales sostenidas, de las que se hace un análisis de correlación pero no se elimina ninguna.
	
	 Para la construcción de modelos se utilizan 4 técnicas diferentes: 3 de ellas de regresión lineal (LS, IRLS y LASSO) y una de regresión no lineal (CART's). Se llega a la conclusión de que los métodos lineales no dan malos resultados, siendo el IRLS el mejor de los 3. Sin embargo, el que mejor precisión da es el método CART. Los errores son de 8.47 $\pm$ 0.17 para UPDRS total con IRLS y de 7.52 $\pm$ 0.25 usando el método CART. A parte de estos resultados, en el artículo se realiza un análisis de la correlación de las características fijándose en los coeficientes devueltos por el método LS. En ellos se puede ver como las características altamente correlacionadas tienen magnitud similar y signo opuesto. 
	
	\item \textit{Novel Speech Signal Processing Algorithms for High-Accuracy Classification of Parkinson's Disease} \cite{MxLtNovel}.\\
	En este artículo el objetivo es clasificar entre personas sanas y con PD a partir de audios de vocales sostenidas \textbf{utilizando un gran conjunto de medidas de disfonía extraídas de los audios}. Un punto novedoso de este experimento, es que hasta ese momento, se habían elegido conjuntos de pocas (<20) medidas de disfonía y medido su correlación. En este artículo, trata un total de 132 medidas lineares y no lineares de cada audio a las que se aplicará diferentes algoritmos de selección de características (algoritmos FS) para elegir varios conjuntos (uno por cada algoritmo). Los algoritmos de selección de características han sido: LASSO, RELIEF, mRMR y LLBFS. Cada algoritmo de selección de características ha elegido un conjunto diferente, siendo básicamente medidas como \english{jitter} y \english{shimmer}, variantes de medidas de ruido, MFCCs y medidas no lineales. Se ha analizado el número de características que comprende cada conjunto y se ha llegado a la conclusión de que cada conjunto tendrá 10 medidas únicamente. Según se explica, esto es debido a que usando más de 22 características se tiende al sobreajuste y que tampoco hay mucha mejoría usando 22 características en lugar de 10.
	
	Posteriormente cada uno de esos conjunto de características ha sido usado para construir dos clasificadores con dos métodos distintos: random forest y SVM con kernel Gaussiano.Se obtienen resultados que mejoran cualquier artículo de de accurracy del 97.7\% utilizando las 132 características con SVM. Sin embargo y como acabamos de explicar anteriormente, en el mismo artículo acusa del accuracy tan alto a un posible sobreajuste al usar tantas caractarísticas. Por lo explicado anteriormente, cada algoritmo FS escoge un conjunto de 10 medidas elegidas. El mejor resultado se da para el subconjunto de características escogido por el algoritmo RELIEF y usando SVM, presentando una accuracy del 98.6\% frente al RF que con el mismo conjunto llega a un accuracy del 93.5\%. 
	
	Para terminar con el resumen de este artículo, indicar que la importancia está en que se trabaja con un número grande de medidas (132), llegando a la conclusión de que los subconjuntos idóneos tendrán alrededor de 10 características ya que tienen resultados parecidos a los conjuntos grandes pero generalizarán mejor que éstos.
	
\end{itemize}

El segundo grupo de investigación y para nosotros más importante, ya que seguimos los pasos para la extracción de características, es el liderado por J. R. Orozco-Arroyave de la Universidad de Antioquía, Colombia. Estos artículos son más recientes que los de el anterior grupo de investigación y, aunque tienen menos repercusión, son autores que trabajan y colaboran con el MIT (Massachusetts Institute of Technology) el cual tiene gran prestigio. El artículo más importante del grupo y que aporta una nueva visión a esta materia es \cite{Orz2016}. Esta nueva visión es debido a que toma en cuenta más audios aparte de únicamente de la vocal sostenida. También tiene en cuenta varios idiomas. Es precisamente de este artículo del que hemos querido seguir los pasos a la hora de determinar qué características sacar de los audios y cómo sacarlas. También valorar y agradecer a este grupo de investigación la cesión de los audios descritos en \cite{OrzCorpus}, un corpus de audios de vocales sostenidas, monólogos, frases y palabras recopilados de un total de 100 personas (50 sanas y 50 PD) que contiene aproximadaamente 4200 audios. Analizamos el artículo más importante de este grupo:
\begin{itemize}
	\item \textit{Automatic detection of Parkinson’s disease in running speech spoken in three different languages} \cite{Orz2016}.\\
	En este artículo clasifica entre personas sanas y con PD y se incorporan a la materia varias novedades relacionadas con los tipos de audios a utilizar.La primera de ellas es utilizar audios de monólogos, texto leído y palabras aparte únicamente de la vocal sostenida. Esto es debido a que la mayoría de los experimentos se basan en vocales sostenidas, obviando los consejos de los neurólogos respecto a que las consonantes requieren un mejor control de los movimientos de órganos (la lengua por ejemplo) y de ahí se puede sacar mucha más información. Por ello se utiliza un método para la caracterización de las señales de voz, basado en la segmentación automática de las expresiones en cuadros con voz y sin voz.
	
	La otra gran novedad que incorpora es la utilización de corpus de audios de varios idiomas a la vez, en este caso 3 diferentes. Esto se utiliza, por ejemplo, para lo que ellos mismos llaman \english{cross-language experiments}. En estos experimentos, el entrenamiento del modelo se hace en un idioma y el testeo con otro idioma diferente. En este artículo se realizan diferentes experimentos, construyendo varios modelos para cada tipo de audio y así evaluar qué tipo de audio es el mejor para la clasificación. Se comienza haciendo un preprocesado de los audios donde los silencios iniciales y finales son eliminados. Se prosigue dividiendo los audios en partes con voz y partes sin voz (estos frames con voz y sin voz se corresponden a los frames donde el software Praat \cite{praat} detecta que hay discurso y donde no lo hace). Para acabar el preprocesado se eliminan los fragmentos menores a 40 ms.
	
	A continuación empiezan con la extracción de características de los audios. Se sacan 3 grupos de medidas diferentes con las que se construyen 3 modelos diferentes: 2 sacando diferentes tipos de medidas de los fragmentos con voz y 1 sacando medidas de los fragmentos sin voz. De los fragmentos con voz, se sacan por un lado 64 medidas prosódicas (\english{jitter}, \english{shimmer}...) con sus estadísticos (media, desviación, coeficiente de asimetría y curtosis). Por otro se sacan NHR, NNE, GNE, \textit{F1}, \textit{F2} y las 12 MFCC con sus estadísticos. De los frames sin voz se hace un análisis cepstral y de energía para calcular las 12 MFCC y 25 BBE con sus estadísticos. 
	
	Después se hace la clasificación utilizando SVM con kernel Gaussiano, que utiliza \english{grid search} para los parámetros \textit{C} que mide la estrictez del margen y $\gamma$ que es un parámetro del kernel Gaussiano utilizando como medida objetivo \english{accuracy}. Se utiliza una técnica de validación cruzada 10-fold.
	
	Como resultado obtenemos 3 modelos diferentes para cada tipo de audio. De cada modelo saca 4 medidas para poderles comparar: \english{accuracy}, \english{sensitivity}, \english{specificity} y \textit{AUC} (área bajo la curva de ROC). En el artículo compara los experimentos por tipo de audio. El mejor rendimiento lo consiguen, el texto leído, ya que consigue, para el idioma castellano y con el modelo construido con las medidas extraídas de los frames sin voz, un accuracy del 97\% y un AUC del 99\%.
	

\end{itemize}

Ahora trataremos dos trabajos de una base da datos adquirida. Estos trabajos corresponden a Giovanni Dimauro de la Universidad de Bari, Italia. En el segundo de estos trabajos se aborda una perspectiva diferente para la diagnosis de Parkinson, utilizando sistemas Speech-To-Text:
\begin{itemize}
	\item \textit{VoxTester, software for digital evaluation of speech changes in Parkinson disease} \cite{giovanni1}.\\
	Este artículo describe la realización de una herramienta software para la ayuda en la evaluación de los cambios y variaciones en la voz de pacientes con la enfermedad del Parkinson. El funcionamiento de la herramienta es el siguiente: insertamos una serie de audios del paciente, la herramienta analiza el audio y nos devuelve 4 gráficas que un experto deberá interpretar para la evaluación de la enfermedad. Las gráficas de salida son las siguientes: gráfica del DDK \english{rate}, gráfica de la intensidad vocal y duración del discurso, gráfica del espectro de frecuencias vocal y gráfica del nivel de presión vocal. Estos parámetros sacados de la voz son importantes ya que pueden ser indicadores de la enfermedad del Parkinson. Por ejemplo, un rango de frecuencias de la voz pequeño, una frecuencia fundamental baja o un decaimiento notable de la intensidad de la voz en el discurso serán una de las características del discurso para personas con esta enfermedad.
	
	En este experimento no se utiliza en ningún momento técnicas de clasificación ni tampoco la herramienta discierne por ella misma si la persona del audio tiene una enfermedad o en qué grado. Lo único que hace es analizar la onda sonora y mostrar algunas características, por lo que no se alinea del todo con nuestros objetivos.
	\item \textit{Assessment of Speech Intelligibility in Parkinson’s Disease Using a Speech-To-Text System} \cite{giovanni2}.\\
	Este experimento trata de resolver el problema del anterior (valuación de los cambios y variaciones en la voz de pacientes con la enfermedad del Parkinson) con otra línea. Para ello utilizará sistemas de reconocimiento de voz y transcripción a texto (\english{Speech-To-Text Systems}) para evaluar la inteligibilidad del discurso de las personas con Parkinson. El proceso será el siguiente: la persona con Parkinson habla a un sistema STT, el texto de la transcripción se pasa a un programa que evalúa su acierto y finalmente ese programa devuelve el porcentaje de fallo en el reconocimiento de voz respecto a la frase objetivo. El objetivo de este proyecto ha sido comparar los fallos de reconocimiento en las palabras entre 3 grupos de personas: jóvenes sanos, mayores sanos y pacientes de la enfermedad del Parkinson. La conclusión a la que llegan los investigadores es que había mucho mas fallo en el reconocimiento de palabras en los pacientes de la enfermedad del Parkinson.
	
\end{itemize}

\tablaSmall{Objetivo de cada artículo}{l c c c}{objetivoporarticulo}
{ \multicolumn{1}{l}{Artículos} & Clasificación PD/HC & Regresión Mapeo UPDRS & Otro análisis\\}{ 
\cite{MxLtSuitability} \cite{MxLtNovel} \cite{Orz2016} & X & &\\
\cite{MxLtAccurate} & & X &\\
\cite{giovanni1} \cite{giovanni2} & & & X\\
} 

\capitulo{7}{Conclusiones y Líneas de trabajo futuras}

\section{Conclusión}
Respecto a los requisitos iniciales del proyecto, pienso que se han cumplido en su mayor medida. Dado a que este proyecto ha sido un proyecto de investigación, hay determinados objetivos que han podido ir cambiando a lo largo del proyecto. Por ello, los objetivos principales de realizar una investigación extensa y documentada y, finalmente, realizar una aplicación base de ejemplo, se han cumplido. 

Cabe destacar que aunque evidentemente todo el proceso de minería de datos y realización de experimentos fue arduo, los primeros pasos de la investigación les considero como los más complejos. Esto es debido a que se quería construir una buena base bien documentada de la investigación y así, poder dar los pasos de manera más firme. Es por eso que las tareas de resumen del estado del arte, identificación de metodologías y adquisición de un conjunto de audios con el que poder trabajar fueron muy costosos por diversos motivos (dificultad de encontrar metodologías, sin respuesta de dueños de las bases de datos...).

En lineas generales, en este trabajo se ha desarrollado una serie de experimentos sobre el tema relacionado de análisis de audios para la detección del Parkinson bastante extenso. Hemos obtenido muchos conjuntos de características de los audios (108 diferentes, extraídos de diferentes formas) con los que se han hecho muchos experimentos. Cada paso que hemos ido dando en el proyecto, ha sido con el objetivo de ir mejorando los resultados obtenidos por nosotros. 

Cabe destacar que los resultados por si solos no son malos, incluso están en la medida de los resultados de algunos artículos (tabla de resultados del estado del arte recogida en \cite{MxLtNovel}). Sin embargo, al estar siguiendo en este proyecto el proceso de \cite{Orz2016}, con el mismo conjunto de audios y las mismas herramientas, cualquier resultado que estuviera lejos del 0,98 o 0,99 AUC nos ha parecido 'bajo'. Destacamos que la base de datos es pequeña ya que, aunque haya 2000 audios, como ya se ha comentado no se pueden entrenar modelos con todos los audios a la vez, hay que dividirlos por tipos, lo que hace que tengamos conjuntos de datos de 50 a 300 instancias.

Si nos olvidamos de la diferencia de resultados con ese artículo concreto, ha sido una labor de experimentación en la labor de Minería de Datos muy satisfactoria. Se han comprendido todas las etapas de la misma, desde la recopilación de datos y la extracción de características, hasta las labores finales de interpretar los resultados. La gran cantidad de diversos experimentos y técnicas utilizadas han servido para aprender en profundidad cómo funcionan los diferentes algoritmos y herramientas de este campo.

En resumen, tanto con la labor de investigación y su documentación, como con la labor de realización de experimentos, como con la realización aplicación final estoy satisfecho.



\section{Líneas futuras}
Se sugieren varios caminos diferentes de investigación, así como varios detalles a mejorar y desarrollar del proyecto:
\begin{itemize}
\item Primeramente, se sugiere que, en vez de realizar una tarea de clasificación entre PD y HC, se realicen \textbf{experimentos de regresión sobre la escala UPDRS}. Es decir, en vez de únicamente decir si una persona tiene Parkinson o no, realizar una regresión para ver que nivel de Parkinson tiene en la escala UPDRS. Abordar el problema de esta manera, se empezó a ver en \cite{MxLtAccurate}. Para esta tarea de regresión se podrá experimentar con todos los conjuntos los conjuntos de características extraídos, incluidos los \english{embeddings}. En principio, iba a ser también uno de los objetivos del proyecto. Sin embargo, al no llegar al rendimiento esperado en clasificación, se optó por realizar más experimentos de clasificación en vez de comenzar con regresión.
\item Con objetivo de mejorar los resultados de los experimentos, se podrán probar diferentes elementos que mejoren su resultado. Uno de ellos es \textbf{probar otros selectores de atributos}. En este proyecto hemos probado PCA, \english{Select K Best} y \english{Variance Treshold}. PCA es un reductor de dimensionalidad. Los otros dos son algoritmos de selección de atributos que se basan en cada atributo individual para su elección. Por ello, como mejora se podrían elegir nuevos selectores de atributos, basados en modelos de aprendizaje (i.e. \english{SelectFromModel} de \textit{Scikit-Learn}) o basados en grupos de atributos (i.e. \english{Recursive feature elimination} de \textit{Scikit-Learn}).
\item Otra perspectiva diferente puede ser \textbf{tratar de manera diferente los \english{embeddings}}. Como hemos comentado, VGGish saca un vector de características para cada fragmento de 0.975s de cada audio. Nuestra solución ha sido hacer la media y desviación para obtener un único vector de características de todos los fragmentos de cada audio. Sin embargo se podrán probar otras perspectivas.
	\begin{itemize}
	\item La primera que se sugiere es, además de sacar la media y desviación, obtener también la \textbf{curtosis y el coeficiente de asimetría} (ver sección \ref{subsec:medidasdistribucion}). Esto ayudará a no desechar tanta información, y puede ser que ayude a mejorar los resultados.
	\item Otra perspectiva que se sugiere es la siguiente: clasificar cada vector devuelto por VGGish por separado, como si fuera un audio diferente, y devolver para ese audio la predicción mayoritaria de las predicciones de sus segmentos. Es decir, para un audio, clasificar cada 0.975s que analiza VGGish por separado de manera independiente y devolver la predicción mayoritaria.
	\end{itemize}
\item Por otro lado, de VGGish únicamente hemos usado la red de extracción de características. También podemos probar a \textbf{modificar la red de clasificación} del proyecto VGGish para que saque 2 salidas, en vez de las 600 que ahora tiene, y probarla sobre nuestro problema.
\item Una mejora que se puede añadir a los experimentos actuales es la siguiente. En el experimento que imita al realizado en \cite{Orz2016}, se utiliza un SVM con kernel gaussiano. La búsqueda de parámetros definida en dicho artículo, dice que buscan el parámetro $C$ y $\lambda$ en los rangos de [1, 10000]. Sin embargo, en los experimentos de la tercera fase, hemos añadido el valor  \textit{auto} para $\lambda$ , que en \english{Scikit-Learn} se corresponde con 1/número\_atribs. Por ello, también se llega a la conclusión, y se sugiere, que en experimentos futuros se realice la \textbf{búsqueda del parámetro $\lambda$} en el rango [1/1000, 1] $\rightarrow$ (1/1000, ..., 1/100, ..., 1/10, ...,1), en vez en el rango [1, 10000].
\item Para terminar, se sugiere la mejora de la aplicación. A parte de mejoras estéticas, se sugiere el uso de \textbf{diferentes clasificadores en función de los detalles} del paciente de entrada y del audio (vocal, palabra o frase). Actualmente utilizamos un clasificador único. La mejora sería derivar la clasificación a un clasificador u a otro en función del sexo, edad y tipo de audio introducido.
\end{itemize}


\medskip

\bibliographystyle{plain}
\bibliography{bibliografia}

\end{document}
